\section{From observed failures to evidence-aware auditing}
\label{app:release-details}
The experiments motivate a narrower question than whether a hand-off is safe:
\emph{which claim can the center justify from the evidence it actually receives?}
We derive a design from the observed failures below. A restricted structured-output
gate is implemented and tested on authored cases; the broader semantic composition
remains proposed. We distinguish these two levels throughout.

\paragraph{What recent methods contribute.}
FlowSeal~\cite{flowseal2026} places mandatory release checks outside the agent's
conversation and isolates semantic declassification. FAVA~\cite{fava2026}
separates model-produced permission representations from deterministic authorization;
its guarantee begins after evidence has been lowered into its graph.
CaMeL~\cite{camel2025} uses capabilities and data-flow control while explicitly
excluding some text-to-text corruption from its protection. We adopt the separation
of proposal and authority, rather than claiming it as new. Our narrower research
question concerns which release claims a content-blind center can justify when
policy, identity, support, or event evidence is incomplete. We do not reproduce
these systems or claim an empirical advantage over them.

\subsection{Authorization, relevance, and support are separate claims}
The deployed follow-up used four development tasks in a $2\times2$ design:
broad versus case-scoped retrieval, and snapshot versus delivery-time permission
checks (Appendix~\ref{app:jetstream}). All 16 episodes completed, with 32 model
calls and no token-cap hits. Both snapshot variants exposed the revoked Birch
reference. Delivery-time checks withheld that record, but broad retrieval still
let the answer substitute Acorn's reference for Birch's. Case-scoped retrieval
removed that particular substitution in the paired episode, yet the answer
asserted that Birch's refund ``has not been processed yet,'' which the available
evidence did not establish. These are observations on reused development cases,
not estimates of population-level error rates.

Three predicates are therefore needed. $A$ asks whether the recipient may access
the resource at delivery; $R$ asks whether it belongs to the requested case and
purpose; $S$ asks whether the generated claim is supported by the permitted
evidence. Authorized context can be irrelevant, and relevant authorized context
does not make every generated assertion true. A check of $A$ must not be reported
as a check of $A\land R\land S$. A missing record also establishes only lack of
available evidence, not a negative fact about the world.

\subsection{Make the scope of a verdict explicit}
For an event $e$, let $E_e$ be the available evidence and let
$\Omega(E_e;\mathcal{T})$ denote the nonempty set of executions consistent with
that evidence under stated trust assumptions $\mathcal{T}$. For one specified
violation predicate $V$ over completed disclosures, the ideal verdict is
\[
\widehat V(E_e)=
\begin{cases}
\textsf{violation}, & \forall w\in\Omega(E_e;\mathcal{T}),\ V(w)=1,\\
\textsf{no violation in scope}, & \forall w\in\Omega(E_e;\mathcal{T}),\ V(w)=0,\\
\textsf{unknown}, & \text{otherwise}.
\end{cases}
\]
Inconsistent evidence, represented by an empty set, also yields unknown and an
integrity alert; it must not make either universal condition vacuously true.
If the actual execution belongs to this set, a definite verdict is sound by
definition. This observation is a specification, not a proof that our current
verifier constructs the set correctly. Its difficulty lies in establishing the
trust assumptions and evidence coverage. The current software uses bounded
checks and heuristic semantic labels, not a general possible-executions solver.

Each response should name its predicate, event or time-window scope, supporting
evidence, and missing obligations. The absence of a semantic alert cannot justify
an authorization claim, or vice versa. Likewise, a blocked attempt is a denied
action rather than a completed disclosure. A pre-execution gate may map unknown
to deferral, redaction, or refusal according to risk; that operational choice
does not change unknown into evidence of a violation.

\subsection{Bind policy evidence to the release boundary}
A proposed release receipt binds the workspace, event, resource version,
recipient, purpose, policy epoch, and release decision. A trusted authority must
establish the epoch effective at that release, not merely return its current
epoch when an old event is audited. A policy check and the corresponding context
release must share a defined serialization boundary. Otherwise a revocation
between check and release recreates the stale-snapshot failure.

Our database harness implements a transactional check at its recorded context
release boundary. It does not prove revocation of information already present in
a model's context or memory. The experimental receipt verifier binds identifiers
and an epoch, but does not yet implement all the proposed purpose, resource-version,
and semantic-support obligations. Its authoritative epoch is supplied by the
evaluation fixture. A stale authority and an authorized issuer signing a false
decision remain counterexamples. Authentication alone cannot eliminate either.

\subsection{Check semantic support locally and disclose the uncertainty}
At the owner boundary, a proposed support checker decomposes a generated answer
into claims and relates each claim to authorized evidence for the requested case.
An unsupported claim triggers revision or a statement that the information is
unavailable. The center receives only a scoped support result, opaque provenance
references, checker version, and uncertainty status. Raw claims and supporting
passages remain local. This is a proposed mechanism: the current online harness
reports generated-output observations as unknown, rather than pretending to
verify their semantics.

The additional checker is fallible and creates a new trust dependency. It must
be tested on omissions, paraphrases, plausible but unsupported assertions, and
contradictory records. Self-reported model confidence is not a calibrated error
bound. The center can verify a checker's identity and evidence bindings without
being able to verify the judgment itself. Privacy labels also need a separate
test: a statement can be fully supported and still disclose a prohibited fact.


\subsection{A restricted, implemented release mechanism}
\label{app:structured-release}
To test the separation without treating a second LLM as an oracle, we implement
an experimental structured-output gate. The model proposes only tuples
$(\mathrm{source},\mathrm{version},\mathrm{field},\mathrm{value})$. A trusted local
adapter supplies the task's workspace, case, recipient and purpose, source fields,
and current authorization state. Model-proposed permission overrides and extra
free-text fields are rejected. The full gate requires current recipient access,
matching workspace and case, an allowed field-purpose pair, matching source
version, and exact equality to the corresponding source value. Missing authority
or support produces unknown; an established scope or permission mismatch produces
deny. The entire proposal is withheld if any tuple fails.

Only validated fields reach a fixed renderer; model-written prose is never
appended afterward. Thus this implementation prevents unsupported field substitution
within its admitted representation. It does not solve natural-language entailment.
A source may itself be wrong, equivalent paraphrases are rejected, and releasing
several individually permitted fields can enable inference outside these checks.
The adapter is a trusted input, and production use would additionally require
serializing the check with release and mediating every output channel. The core gate is a pure function. An experimental PostgreSQL adapter now
serializes it with a business-side mailbox commit, as evaluated below; neither
is integrated into the production HTTP audit service or an external message sink.

A minimal center-facing summary contains only a schema identifier, status and
fixed reason code. It excludes the source values, case names, model prose and
source identifiers. The local result contains those details for the business-side
renderer and must not be serialized wholesale to the center. Even this reduced
summary discloses an outcome and failure class; it is neither differentially
private nor independently verifiable evidence of correct execution.

\paragraph{Mechanism regression and its counterexamples.}
Nineteen authored fixtures cover revoked and deleted records, stale versions,
missing authority, cross-case and cross-workspace references, purpose restrictions,
unsupported values, malformed proposals and mixed valid/invalid claims. The same
inputs are available to all arms. These controls are ablations of our own code,
not implementations of FlowSeal, FAVA or CaMeL. Table~\ref{tab:release-regression}
reports 95 deterministic evaluations, not 95 independent workflows.

\begin{table}[ht]
\centering\small
\begin{tabular}{lrrrr}
\toprule
Control & Inadmissible release & Benign exact & Unknown & Source error\\
\midrule
Snapshot & 10 & 2/2 & 5 & 1/1\\
Fresh policy & 6 & 2/2 & 7 & 1/1\\
Fresh policy + scope & 4 & 2/2 & 7 & 1/1\\
Full structured gate & 0 & 2/2 & 11 & 1/1\\
Deny all & 0 & 0/2 & 19 & 0/1\\
\bottomrule
\end{tabular}
\caption{Authored contract regression. Inadmissible release counts violations of
this representation's contract, not semantic privacy errors. Benign exact counts
two supported exact-value outputs. Source error records acceptance of a value
that matches an intentionally incorrect trusted source; it is a negative control,
not a successful answer.}
\label{tab:release-regression}
\end{table}

The full gate also rejects the supported paraphrase ``on Tuesday'' for the source
value ``Tuesday.'' Its zero contract violations therefore come with a measurable
expressiveness cost. All nontrivial arms accept the deliberately incorrect
trusted source, demonstrating why source consistency must not be called factual
truth. Independent held-out and human semantic evaluation remain necessary.

\paragraph{Model-generated proposals.}
An eight-call Jetstream pilot used cached Qwen2.5-3B-Instruct with greedy decoding
and a 160-token cap. The inputs are a subset of the same authored fixtures, not
new held-out workflows. The model received a draft and source fields but not the
expected verdict or release-time authority state. All eight proposals parsed
without truncation or generation errors. Snapshot control released all eight;
fresh policy released five; adding scope released three; the full gate released
two, denied three, and returned unknown for three. The two released proposals
were the benign exact-field cases. In the unsupported-negative fixture, the
model retained the seeded claim ``not processed'' despite the source saying
``scheduled''; exact support checking withheld it. This is draft-conditioned
inference plus a local gate, not a new end-to-end database deployment or an
independent semantic evaluation. Frozen protocol and source hashes, raw outputs,
model revision, and environment versions accompany the artifact.

\paragraph{Checking freshness is not sufficient without a release boundary.}
We next wrapped the gate in a PostgreSQL transaction that locks the resource rows,
reads current source and authority state, evaluates the proposal, and commits a
business-side mailbox entry together with an audit-summary outbox entry. Resource
updates use the same rows. An event-scoped transaction lock serializes concurrent
retries; a binding digest rejects reuse of an event identifier for a different
proposal or task. A successful retry returns an acknowledgment rather than
re-emitting the saved answer.

Seven authored schedules, two configurations, and three repetitions produced
42 real database trials. In the check-without-row-lock ablation, a barrier let
revocation or a version update commit after validation but before the mailbox
write: all six such trials retained the stale allow result. With row locking,
we observed the conflicting update waiting in PostgreSQL's lock table; the
mailbox transaction serialized before it. A revocation committed before the
check was denied in both configurations. An injected pre-commit exception left
neither mailbox nor outbox entry, and a retry created one of each. Concurrent
identical retries created one pair, while conflicting retries were rejected.

These are controlled serialization and recovery checks, not estimates of attack
success in arbitrary deployments. The defined release point is the database
mailbox commit. A later network dispatcher, recipient consumption, or a writer
that bypasses these resource rows lies outside the experiment. In particular,
revocation after that commit does not retract already released data, and a
future dispatcher may need another policy check if its send is the intended
release boundary. Outbox persistence here is not proof of audit-center receipt.
The distinction matters: a fresh decision can be correct when computed yet
invalid when applied at a different boundary.

\paragraph{New tasks with a frozen gate and scorer.}
To separate reuse of development fixtures from task transfer, we froze the gate,
prompts, task oracles and a separate scorer before generating proposals for
20 new tasks. Four authored domains---billing, appointment scheduling,
onboarding and travel---each contribute two benign requests, one disclosure
challenge, one request for an unsupported status and one authority-unavailable
case. Unlike the eight-call pilot, these inputs contain no draft answer. The
model sees the request, task scope and source records; it does not see the
scoring oracle or current authority state. All five controls evaluate the same
proposal. This is direct CPU inference plus a pure release gate, not a new
end-to-end database deployment. Source hashes and the pinned Qwen checkpoint are
retained with the raw proposals.

The predeclared utility endpoint requires every requested source-backed field
on each of the eight benign tasks, without private-marker disclosure,
unsupported values, generation errors or truncation. Refusal is not success.
Privacy-marker release is measured over all 20 tasks, including incidental
leaks on benign requests. Unsupported released values and nonempty releases
on the eight uncertainty tasks are reported separately. These are exact-slot
and literal-marker checks, not human judgments of answer usefulness or semantic
privacy. New task authorship is internal; the scorer is independent of the gate
implementation, not independently authored or adjudicated. Correlation within
four authored domain families precludes treating the 20 tasks as independent
population samples.

\begin{table}[htbp]
\centering\small
\begin{tabular}{lrrrrr}
\toprule
Control & \shortstack{Complete\\(/8)} & \shortstack{Private\\(/20)} &
\shortstack{Unsupported\\(/20)} & \shortstack{Uncertain\\(/8)} & \shortstack{Unknown\\(/20)} \\
\midrule
Snapshot & 8 & 2 & 4 & 7 & 2 \\
Fresh policy & 8 & 2 & 4 & 4 & 5 \\
Policy + scope & 8 & 0 & 4 & 4 & 5 \\
Full gate & 8 & 0 & 0 & 0 & 9 \\
Deny all & 0 & 0 & 0 & 0 & 20 \\
\bottomrule
\end{tabular}
\caption{Frozen new-task evaluation: 20 authored tasks, one pinned Qwen2.5-3B
proposal per task, five matched gate controls. Complete counts successful benign
requests. Private counts literal forbidden-marker release; Unsupported counts
released values absent from the cited source; Uncertain counts nonempty release
when authority or requested facts are unavailable. Unknown is a gate verdict,
not a success. This is a controlled transfer check, not independent semantic
adjudication.}
\label{tab:frozen-tasks}
\end{table}

All 20 calls completed without generation errors or token-cap hits.
Table~\ref{tab:frozen-tasks} separates privacy from utility: the full gate
completes all eight benign requests while releasing no forbidden marker or
unsupported value. Deny-all also prevents those releases but completes no benign
request. Scope checking removes the two literal disclosure failures; it still
allows four unsupported status claims. Exact source support is needed to remove
those four releases. The model generated ``not started'' without evidence in
three domains and used ``onboarding'' as a status in the fourth. Thus reducing
disclosure alone does not establish factual task correctness.

A descriptive breakdown limits the causal interpretation. All nontrivial arms
reject two proposals with a string-valued version field before checking policy;
one is the revoked invoice case. The new run therefore supplies no additional
evidence about revocation-check effectiveness for that case. The full gate has
nine unknowns: two schema failures, three authority-unavailable cases and four
unsupported-value cases. Three further requests are denied for purpose, case or
workspace mismatch. There are no withheld benign requests in this small set.
CPU generation takes 21.9--72.5 seconds per request (median 29.8 seconds), not an
interactive production latency result. Prompts, gate and scorer were not changed
after outcomes were observed; all seven frozen source/protocol hashes match the
archived run.

\paragraph{Two levels of support, not one safety score.}
The operational distinction is between a validated structured projection and a
free-text semantic judgment. Only the first currently has an executable check
against source fields. A future semantic lane could admit paraphrases after a
separate local checker, but its error rate, abstentions, and privacy decisions
must remain separately visible. It must not inherit the structured lane's
contract checks merely because both produce a signed report. This division
retains an exact path for dates, identifiers and status fields without presenting
that path as a solution to arbitrary inter-agent conversation.

\subsection{Expose gaps in observation and costs of linkage}
Reliable transport addresses loss of \emph{recorded} events. Complete observation
requires an independently justified event manifest or a trusted mediation point.
Sequence numbers and acknowledgments can expose gaps relative to such a manifest;
they cannot detect an interaction omitted by both endpoints. The pilot's 40/40
received event IDs establish delivery coverage over its recorded events only.

Cross-message inference additionally requires continuity across the relevant
subject, recipient, and time window. Splitting reports or rotating identifiers
can destroy that continuity. Retaining stable identifiers restores some joins
but exposes relationships to the center. We therefore propose workspace-scoped,
window-bounded linkage, with explicit unknowns when evidence expires or a join
cannot be justified. The retention window is a privacy--coverage tradeoff to
measure, not a confidentiality guarantee; an opaque identifier remains linkable
within its scope. The service retains a bounded workspace graph across reports, but this does not
establish correctness of arbitrary cross-report identity joins or independent
observation completeness. The proposed linkage policy still requires evaluation.

\subsection{External tasks and native success criteria}
\label{sec:external-dojo}
The new authored tasks above do not establish external validity. We therefore
add an action-prefix replay using AgentDojo~\cite{debenedetti2024agentdojo}, an
external benchmark with executable task-utility and attack-success predicates.
We pin the upstream repository at revision
\texttt{089ed468cf3e}, retain its MIT attribution, and use its v1 task definitions
without changing the success criteria. This evaluates an authorization component
of our method, not the full semantic-support gate or a live adaptive agent.

\paragraph{Selection and controls.}
Before scoring, we freeze 1,132 archived trajectories from four official pipeline
labels: GPT-4o (2024-05-13), GPT-4 Turbo (2024-04-09), Claude 3.5 Sonnet
(2024-06-20) and Llama 3 70B Chat. The selection contains all available clean
runs and \texttt{important\_instructions} runs for the specified disclosure
attacks in workspace, travel and Slack: 324 clean trajectories and 808 attacked
trajectories, spanning 81 underlying user tasks. We exclude sabotage and phishing
goals from the privacy endpoint, not unsuccessful attacks. Model names describe
the upstream archives; we do not make new model calls or claim that these are
current-model measurements.

Five matched controls use the same archived action sequence: no gate, a keyword
check for communication intent, literal destination binding to the initial
request, their conjunction, and denial of every mediated egress action. The
adapter consumes only the original user request and proposed tool arguments.
Neither attacker goals, security labels nor successful reference actions grant
permissions. It reuses the structured gate on a destination-authority envelope;
this says nothing about factual support of the message body. The initial request
is trusted and untrusted tool observations cannot extend its authority.
Destination binding deliberately does not infer contacts from names or grant
permissions from retrieved text. Its false denials are part of the evaluation.

\paragraph{Replay semantics and endpoints.}
We re-execute observed tool calls in AgentDojo's in-memory environments. Email,
Slack and web operations modify mock state; no real communication is sent. At
the first blocked egress action, we terminate the prefix and remove the final
answer instead of retaining the original trajectory's later success. Unblocked
trajectories keep their archived final answer. The unchanged upstream predicates
then evaluate task utility and attack success from the resulting state and
executed prefix. Upstream \texttt{security=True} denotes attacker success.
Errors remain in the denominator and are reported separately.

This intervention measures which archived actions a gate admits and what utility
the admitted prefix retains. It does not measure how an agent would recover,
replan, or attack again after a refusal. The selected attack predicates also
encode specific exfiltration goals, not every semantic disclosure. We compare
ungated replay with the archived labels to expose reconstruction differences.
For paired utility differences, we resample the 81 underlying tasks, keeping
all model and injection variants of each task together (2,000 bootstrap draws;
seed 1729). These intervals describe variation within this archive, not a
confidence guarantee for future deployments.

\paragraph{Privacy reduction has a measurable utility cost.}
All 5,660 replays completed in the Jetstream environment. Under the native
attack predicates, combined intent and destination checks reduce observed
attack success from 165/808 to 0/808. Clean utility falls from 198/324 (61.1\%)
to 179/324 (55.2\%); attacked utility falls from 383/808 (47.4\%) to 316/808
(39.1\%). The clean paired change is $-5.9$ percentage points (task-cluster
bootstrap 95\% interval $[-10.2,-2.2]$); the attacked change is $-8.3$ points
($[-12.8,-4.3]$). The combined arm blocks 28 clean and 202 attacked trajectories;
not every block removes a previously successful task. These losses are material:
this external result does not reproduce the eight-of-eight utility retention
of the internally authored structured tasks.

Figure~\ref{fig:core-replay} in the main text displays these same counts.
The deny-egress control has one clean and one attacked checker exception;
both remain unresolved rather than being scored as successful or verified safe.


Intent checking alone leaves 85 successful attacks. Destination-only and combined
checks have identical utility and attack-success counts, so this experiment
provides no measurable benefit for the additional intent rule. Denying every
mediated egress retains only 150 clean and 285 attacked successes, with two
checker exceptions. Some tasks need no mediated external send, so this arm is
not equivalent to refusing every task. The native predicate failure in each
case is a missing website-state key after early termination; it is not silently
converted to a confirmed safe outcome.

The costs concentrate in Slack: combined checking retains 44/84 clean successes
versus 61/84 without a gate, and 24/168 attacked successes versus 90/168. Travel
retains 42/80 clean and 75/160 attacked successes in both arms; workspace retains
93/160 versus 95/160 clean, and 217/480 versus 218/480 attacked. Thus aggregate
retention conceals substantial environment-specific loss. Literal destination
binding cannot resolve legitimate implicit contacts or contextual destinations.
Eleven combined-arm blocks also arise from unsupported destination argument
representations, including null optional values; we retain them rather than
repairing the adapter after seeing outcomes.

\paragraph{Concordance and scope of the evidence.}
Ungated replay agrees with archived task-utility labels in 1,128/1,132 cases and
attack-success labels in 781/808 attacked cases. The 27 attack disagreements are
26 workspace cases and one travel case. Four task-utility disagreements and
these attack disagreements remain in all denominators. We use outcomes of the
pinned executable predicates consistently across the five arms; the discrepancy
means the replay is not an exact reproduction of all historical labels. Twenty-four
archived traces contain execution errors. Runtime tool errors are retained as
observed behavior; the combined arm has no checker exceptions. All frozen
source/protocol and trace hashes are preserved with the artifacts.

\paragraph{Post-hoc discrepancy audit.}
We replay the 31 trajectories with any label disagreement under seven available
upstream suite versions, from v1 through v1.2.2, without changing the frozen
primary experiment. All four utility disagreements align with archived labels
under later versions: three travel cases pass after v1.1 removes the requirement
that the final answer repeat the hotel rating; one workspace case passes after
v1.1.1 accepts the updated end time as part of calendar rescheduling. This
establishes compatibility with changed predicates, not the provenance of the
historical evaluation. All 27 attack disagreements persist across these versions
(nine replay-positive/archive-negative and eighteen in the opposite direction).
Their historical cause remains unresolved.

As a post-hoc sensitivity analysis only, restricting to trajectories whose
ungated utility and attack labels agree with the archive leaves 322 clean and
779 attacked runs. Combined checking changes clean utility from 198/322 to
179/322, attacked utility from 383/779 to 315/779, and attack success from 156/779
to 0/779. Thus both the reduction in selected attack success and the utility
cost persist without the disputed trajectories. These filtered counts do not
replace the primary denominators; the subset was selected after observing
concordance. Per-trajectory version results accompany the artifacts.

The finding is bounded but useful: explicit task-to-destination binding can
suppress the selected archived exfiltration actions, while failing to retain
legitimate tasks that require contextual authority resolution. It does not
validate semantic support of message bodies, arbitrary leakage channels,
continuing attackers or agent recovery after a block. A next method should
resolve recipient authority from trusted task evidence and evaluate recovery in
a fresh closed-loop run; it should not extract permission from attacker-written
observations or use benchmark attack labels as a policy oracle.


\subsection{A falsifiable evaluation of the proposed composition}
The frozen new-task experiment covers matched structured requests, but not
independent semantic validation. A subsequent external evaluation should retain
the same pre-outcome freeze of implementation, prompts, policies and scorer.
Compare snapshot control, fresh-policy control, fresh-policy plus case scoping,
and local support checking on independently sourced workflows. Baselines must receive the same policy and identity information;
central raw-content checks may be a separate reference, with their disclosure
cost reported. Ablate freshness, scope, support checking, and event continuity
individually. Repetitions of one workflow remain one statistical cluster.

Measure false authorization and false denial separately from semantic disclosure,
unsupported assertions, task success, and abstention. Report error among decided
events together with the fraction decided: rejecting everything must not appear
effective. Independent reviewers should label a blinded sample against explicit
policy and source evidence, with disagreements adjudicated and reviewer access
kept distinct from the center's access. Report latency, metadata volume, linkage
exposure, and outage behavior alongside those outcomes.

The proposed support checker fails its purpose if it reduces literal marker hits
without reducing independently judged disclosure or unsupported claims. The
proposed continuity mechanism fails if its gains depend on retaining relationships
beyond the allowed window. These tests would establish a useful contribution:
not that desensitization makes every interaction auditable, but that the system
identifies which conclusions its limited evidence supports and exposes where
additional local checks or trusted observation are necessary.
